% =====================================================================
\chapter{Conclusiones}
\label{cap:conclusiones}

\section{Respuesta a los objetivos}

\begin{enumerate}
  \item \textbf{Corpus de noticias.} Se construyó un corpus de 14.012 notas de siete portales
        argentinos, auditado y depurado de duplicados, publicaciones cruzadas, plantillas y
        contenido que no era noticia, y estructurado como serie temporal por portal, fecha y
        rueda bursátil. La regla de las 17:00 asignó cada nota a la primera rueda en la que el
        mercado pudo reaccionar a ella.
  \item \textbf{Modelo de sentimiento.} FinBETO, un BETO con preentrenamiento continuado sobre el
        corpus y ajuste fino en cuatro clases, alcanzó un F1 macro de 0,685 sobre noticias no
        vistas, y de 0,724 frente a 0,546 del modelo de referencia en tres clases. La medición
        corrigió una fuga de información de la evaluación original, que la sobrestimaba.
  \item \textbf{Trending topics de X/Twitter.} No se realizó: los datos no llegaron a
        integrarse al análisis.
  \item \textbf{Variables de shock.} Se definieron y cuantificaron la atención anormal
        (volumetría de publicaciones), el tono anormal (temperatura del sentimiento), las
        ruedas de tono extremo, el tono propio de cada empresa a partir del reconocimiento de
        entidades, y el shock de cantidad con la definición literal del plan.
  \item \textbf{Agrupamiento de empresas.} El clustering de las 47 acciones no encontró
        familias de comportamiento: el único grupo de tamaño real se explica por la baja
        liquidez de sus integrantes. La sensibilidad al flujo informativo se analizó por
        empresa y por sector.
  \item \textbf{Relación con desfases.} Se modeló con VAR, causalidad de Granger,
        impulso-respuesta, análisis por régimen, una prueba con las notas publicadas con el
        mercado cerrado, dos estudios de eventos y doce variantes de robustez. El DLM previsto
        en el plan no se implementó; el VAR ya incorpora los efectos rezagados.
\end{enumerate}

\section{La hipótesis}

La hipótesis H1 sostenía que el shock informativo tiene un impacto medible y significativo
sobre el precio o el volumen de las acciones argentinas. \textbf{En su forma principal —que el
tono de las noticias anticipe los precios— no se sostiene}: el tono no anticipó el retorno ni el
volumen en el mercado, en las empresas, en los sectores ni en ningún régimen, ni siquiera con la
prueba más favorable, que usa solo las notas publicadas antes de la apertura. Lo que apareció
con claridad es la \textbf{relación inversa}: el retorno del S\&P~Merval anticipa el tono de las
notas siguientes, en 39 de las 47 empresas y en 15 de las 17 con tono propio, y el efecto se
concentra en las notas publicadas la misma noche, después del cierre. Los medios económicos
argentinos, en el período estudiado, contaron lo que el mercado ya había hecho.

Con la definición de shock del plan —muchas más notas que lo habitual— H1 encuentra un
\textbf{apoyo parcial, en el volumen}: cuando una empresa aparece en muchas más notas, ese día
su acción se opera más. Ese mayor volumen no viene con una dirección en el precio, no se
arrastra a los días siguientes y no se anticipa con las notas de la noche. La cantidad de
cobertura acompaña a la actividad del mercado, pero no la anticipa.

\section{Aporte}

El trabajo aporta el primer cruce, hasta donde se pudo relevar, entre el tono de la prensa
económica argentina y el precio y el volumen de las acciones de BYMA; un conjunto de oraciones
de prensa económica argentina etiquetadas por sentimiento; una medida de tono por empresa que
descuenta el clima general; y un \textit{pipeline} reproducible de punta a punta, con controles
automáticos, especificaciones escritas antes de correr los análisis y corrección por pruebas
múltiples.

\section{Limitaciones}
\label{sec:limitaciones}

\begin{itemize}
  \item \textbf{Un año de datos.} Con 224 ruedas en los modelos, un efecto chico de las noticias
        sobre los precios podría existir y no detectarse. El indicio de la prueba pre-apertura
        ($p\approx0{,}1$) es compatible con un efecto chico o con ninguno.
  \item \textbf{Etiquetas de un modelo, no de personas.} El F1 mide el acuerdo de FinBETO con
        Gemma, no con un criterio humano, y no se hizo una validación con anotadores. El modelo
        que clasificó el corpus es el entrenado con la división original; la corrección se
        aplicó a la medición de su calidad, no al modelo.
  \item \textbf{Tono por empresa.} Es el tono de las oraciones que mencionan a la empresa, no una
        opinión sobre ella: una oración negativa sobre el sector que la nombra cuenta como
        negativa para ella. En muchas ruedas se apoya en una o dos oraciones.
  \item \textbf{Hora de publicación.} Es la que informa cada portal. El 13,4\% de las notas tiene
        horas por defecto; los resultados se repitieron sin ellas.
  \item \textbf{Datos diarios.} No permiten ordenar lo que pasa dentro de una misma rueda; la
        prueba pre-apertura es la aproximación más cercana. Algunas acciones argentinas cotizan
        también en Nueva York después del cierre local, y las notas de la noche pueden estar
        contando esos movimientos.
  \item \textbf{Pesos nominales.} Precios, retornos y montos incluyen inflación y tipo de cambio.
        Las variables anormales comparan con las 20 ruedas previas, lo que descuenta la
        tendencia lenta, pero no la elimina del todo.
  \item \textbf{Cobertura del corpus.} Depende de qué capturó la Wayback Machine y de la
        retención de los mapas del sitio; abril de 2026 tiene un hueco de recolección en
        iProfesional, el portal más grande.
  \item \textbf{Volatilidad simétrica.} El VIX argentino sube igual con una suba que con una
        caída, mientras que la evidencia para el S\&P~Merval indica que las caídas pesan más
        \citep{otuki2008}; y hereda el eco de la ventana de 20 ruedas, que también afecta a las
        fechas de los regímenes.
  \item \textbf{Pocos eventos de mercado.} Los estudios de eventos del mercado tienen entre 5 y
        16 eventos, por lo que se apoyan en la comparación contra fechas al azar.
\end{itemize}

\section{Líneas futuras}
\label{sec:futuro}

\begin{itemize}
  \item \textbf{Validar el etiquetado con personas.} Una muestra chica etiquetada por
        anotadores alcanzaría para justificar estadísticamente el uso de Gemma como
        etiquetador \citep{calderon2025alttest}, y permitiría reentrenar FinBETO con los datos
        limpios y una división por fecha \citep{xue2026temporal}. Probar MarIA como modelo base
        es la alternativa natural \citep{gutierrez2022maria,pan2023peerj}.
  \item \textbf{Redes sociales.} Integrar los \textit{trending topics} de X/Twitter, como
        preveía el plan, para validar los shocks que aparecen a la vez en medios y en redes
        \citep{yang2024,wu2026}.
  \item \textbf{Modelo de rezagos distribuidos.} Implementar el DLM del plan como
        complemento del VAR.
  \item \textbf{Análisis en dólares.} Repetir los modelos con precios en dólares contado con
        liquidación, para separar el efecto de las noticias del de la inflación y el tipo de
        cambio.
  \item \textbf{Precios intradiarios.} Con datos hora a hora se podría ordenar noticias y
        precios dentro de la rueda, la principal limitación de este trabajo.
  \item \textbf{Más información por nota.} Separar las noticias por tema \citep{chen2025topics},
        medir cuán central es cada empresa en la nota \citep{asgarieh2024salient} y ampliar el
        corpus con otras fuentes, como GDELT \citep{fronzetti2026gdelt}.
  \item \textbf{Volatilidad.} Reemplazar la media móvil de 20 ruedas por una media exponencial
        \citep{riskmetrics1996}, que olvida los shocks de a poco en lugar de soltarlos de golpe.
  \item \textbf{Ventana más larga.} Con varios años de datos, que incluyan recesiones, se podría
        probar si el tono anticipa al mercado en los períodos de crisis, como encontró
        \citet{garcia2013}.
\end{itemize}
